% =====================================================================
%  2bat-ud1-12.tex  —  SESSIÓ 12: Full de càlcul de gestió
%  (sense preàmbul: s'inclou des de main.tex)
%  Criteris: C2 i C3. Sessió TECNOLÒGICA (full de gestió "viu").
% =====================================================================
\horatitol{Sessió 12 — Full de càlcul de gestió}%
{Muntar un full "viu" que enllaça matrices amb fórmules i recalcula sol en canviar una dada.}

\subsection{Idea clau}

A la sessió anterior vas fer operacions soltes amb el full de càlcul. Ara faràs un
pas més: muntar un \textbf{full de gestió} on totes les dades estiguin
\emph{enllaçades} amb fórmules. L'objectiu és un full \textbf{viu}: quan canvies
una sola dada d'entrada (l'ocupació d'un servei, el percentatge de retallada, les
hores d'una sessió), \emph{tots} els totals, sumes i productes es recalculen sols
a l'instant. Aquesta és la potència real del full de càlcul com a eina de gestió i
de simulació d'escenaris, la mateixa que faràs servir als cicles formatius.

\begin{idea}
\textbf{Idees clau d'un full de gestió}
\begin{itemize}[leftmargin=1.4em, itemsep=2pt]
  \item \textbf{Dades d'entrada} separades dels \textbf{resultats}: les dades es
        posen una sola vegada; els resultats es calculen amb fórmules que hi fan
        referència.
  \item \textbf{Fórmules enllaçades:} cap resultat es tecleja a mà; tot surt de
        fórmules. Així el full es manté coherent.
  \item \textbf{Recàlcul automàtic:} si canvia una dada d'entrada, tot el que en
        depèn s'actualitza sol.
  \item \textbf{Producte de matrices:} moltes fulls de càlcul tenen una funció per
        multiplicar matrices (per exemple \texttt{=MMULT(rang1;rang2)}), però
        també es pot fer cel·la a cel·la amb sumes de productes.
\end{itemize}
\end{idea}

\subsection{Exemples}

\begin{exemple}[un total que es recalcula sol]
Tens l'ocupació de places d'un servei a les cel·les \texttt{B2:B4}. Si a la cel·la
\texttt{B5} escrius la fórmula
\[
\texttt{=B2+B3+B4}
\]
obtens el total. Si després canvies qualsevol valor de \texttt{B2:B4}, el total de
\texttt{B5} es \textbf{recalcula sol}: no cal tornar a sumar res a mà.
\end{exemple}

\begin{exemple}[un producte amb referència a un percentatge]
Si tens un pressupost total al rang \texttt{A1:B2} i el percentatge de retallada a
la cel·la \texttt{D1} (per exemple, $0{,}9$ per a una retallada del $10\%$), pots
escriure a \texttt{A4} la fórmula
\[
\texttt{=A1*\$D\$1}
\]
i arrossegar-la al bloc \texttt{A4:B5}. El símbol \texttt{\$} \emph{fixa} la cel·la
\texttt{D1} perquè no es desplaci en arrossegar. Si canvies \texttt{D1} a $0{,}8$,
tot el pressupost retallat es recalcula sol.
\end{exemple}

% ---------------------------------------------------------------------
\subsection{Guia pràctica (pas a pas)}

Muntaràs un full de gestió de la dedicació professional, reprenent el cas de la
sessió 8: una matriu de sessions per usuari i una columna d'hores per sessió.
\[
S=\begin{pmatrix} 3 & 2 & 1 \\ 1 & 4 & 2 \\ 2 & 1 & 3 \end{pmatrix},\qquad
H=\begin{pmatrix} 2 \\ 1 \\ 1 \end{pmatrix}.
\]

\begin{idea}
\textbf{Passos a l'ordinador}
\begin{enumerate}[leftmargin=1.8em, itemsep=3pt]
  \item Escriu la matriu $S$ al rang \texttt{A1:C3} i la columna d'hores $H$ al
        rang \texttt{E1:E3}.
  \item \textbf{Hores de cada usuari (producte).} A la cel·la \texttt{G1} escriu
        la suma de productes de la fila 1 de $S$ per la columna $H$:
        \texttt{=A1*\$E\$1+B1*\$E\$2+C1*\$E\$3}. Arrossega-la fins a \texttt{G3}
        per obtenir les hores dels tres usuaris. (Si la coneixes, pots usar
        \texttt{=MMULT(A1:C3;E1:E3)}.)
  \item \textbf{Total d'hores.} A la cel·la \texttt{G5} escriu \texttt{=G1+G2+G3}
        per saber les hores de professional que cal preveure en total.
  \item \textbf{Prova un escenari.} Canvia la durada de la primera sessió: posa $3$
        a la cel·la \texttt{E1}. Observa com es recalculen soles totes les hores i
        el total. Anota el nou total.
  \item \textbf{Comprova} que els resultats coincideixen amb el càlcul a mà de la
        sessió 8.
\end{enumerate}
\end{idea}

\textbf{Tasques per lliurar}
\begin{exercicis}
\begin{exlist}
  \item Munta el full descrit i anota les hores de cada usuari ($G1$ a $G3$) i el
        total ($G5$) amb les dades originals.
  \item Prova l'escenari del pas 4 (durada de la primera sessió $=3$) i anota el
        nou total. Explica, en una frase, per què ha canviat sense tornar a
        escriure cap fórmula.
\end{exlist}
\end{exercicis}

% ---------------------------------------------------------------------
\fullsolucions{Sessió 12}
\begin{solucions}
\begin{sollist}
  \item Amb les dades originals ($H=\begin{pmatrix} 2 \\ 1 \\ 1 \end{pmatrix}$), les
        hores de cada usuari són $\begin{pmatrix} 9 \\ 8 \\ 8 \end{pmatrix}$ i el
        total és $9+8+8=25$ hores. Coincideix amb el càlcul a mà de la sessió 8.
  \item Amb la durada de la primera sessió a $3$ ($H=\begin{pmatrix} 3 \\ 1 \\ 1 \end{pmatrix}$),
        les hores passen a $\begin{pmatrix} 12 \\ 9 \\ 10 \end{pmatrix}$ i el total
        a $12+9+10=31$ hores. Ha canviat sol perquè les fórmules fan referència a
        la cel·la \texttt{E1}: en modificar-la, tots els productes i el total es
        recalculen automàticament.
\end{sollist}

\noindent\textbf{Nota per al professorat.} S'avalua la \emph{correcció tècnica}
(matrices en rangs, fórmules de producte amb referències fixes \texttt{\$}, full
enllaçat i coherent) i la \emph{interpretació} (entendre el recàlcul automàtic i
saber llegir l'efecte d'un escenari sobre el conjunt). És la pràctica principal de
la unitat per a l'instrument «Treballs pràctics i tecnològics» (15\%).
\end{solucions}
